\section{Introduction}
\label{sec:intro}

How long does an optimizer stay in a loss basin?
Arguments for flat-minima bias, implicit regularization and learning-rate schedules all depend on the answer \citep{xie2021diffusion,mori2022powerlaw,smith2018bayesian}.
Those arguments lean on mean escape-time formulas derived in asymptotic limits.

The mean is not enough.
The full distribution tells us whether escape is memoryless, which is what the usual stochastic differential equation (SDE) picture assumes.
In the high-barrier limit, Kramers theory predicts Poisson exits and therefore exponential escape times \citep{kramers1940brownian}.
Real optimizers run at finite noise, start at a specific point, and move through landscapes with many neighbouring minima.
None of these conditions guarantees exponential times, and each one bends the distribution in its own way.

Prior work measures mean rates \citep{xie2021diffusion,mori2022powerlaw,ibayashi2021exponential}.
We found no study of distribution shape with censoring and model comparison in real networks, although keyword searches of arXiv are not proof of absence.
Theory also disagrees on direction: \citet{zhou2020towards} predict that SGD escapes a given basin faster than Adam, while a common intuition holds the reverse.
On one hand, the SDE view gives clean scaling laws.
On the other hand, few measurements test those laws against exact minima and saddles.

We test four numerical claims about escape.
(H1) Escape times are Weibull, not exponential, with shape $k$ in 0.7--1.3 depending on geometry.
(H2) Adam escapes narrow basins 2--5$\times$ faster than SGD, with similar times at wide basins.
(H3) Doubling the batch size $B$ at fixed learning rate $\eta$ multiplies mean escape time by $\sqrt{2}$, with $B/\eta$ setting temperature.
(H4) Hessian eigenvalue gaps at saddles predict the escape direction with 70--85\% accuracy.
We fix the pass criteria before any experiment.
We then run hundreds of independent replicas of real SGD and Adam from the same exact minimum or saddle of a small tanh network, and we decide exit by quenching each replica to the minimum it belongs to (see \algref{alg:quench}).

None of the four claims holds as stated, and each fails in an informative way.
Exponential times are rejected in 86\% of conditions, but only 32\% of fitted $k$ values lie in 0.7--1.3 (median 1.47), and $k$ shows no relation to sharpness.
Discarding the earliest quarter of exits moves median $k$ from 1.47 to 1.03.
Doubling $B$ multiplies escape time by 5.1 [3.4, 8.0], about $\tau\propto B^{2}$, not $\sqrt{B}$.
Adam is about $2\times$ faster at matched in-basin loss and slower by at least 7--34$\times$ at its default learning rate.
At low noise, saddle exit follows the top eigenvector in 82--86\% of runs, yet the gap does not predict that share.

\noindent Our contributions are:
\begin{itemize}[leftmargin=*,itemsep=0pt,topsep=0pt]
    \item We build an escape-time instrument for real networks from exact critical points, converged quenching and interval-censored fits, and we validate it against the exact Kramers rate.
    \item We measure the shape of escape-time distributions across 28 well-observed conditions, and we decompose non-exponentiality into a start-up delay and heterogeneity across minima.
    \item We show that batch size acts far more strongly than a $\sqrt{B}$ law, and that the sign of the Adam--SGD comparison depends on learning-rate calibration.
    \item We show that saddle exit is highly reproducible even when it leaves the top eigenvector, which a linearized model cannot explain.
\end{itemize}

\vspace{4pt}
\noindent \Secref{sec:related} reviews prior work, \secref{sec:method} describes our method, \secref{sec:results} reports results, and \secref{sec:discussion} discusses them.
